\section{A Sufficient Condition for a Single Pass}
\label{sec:existence}

归约环节在每个位置 $j$ 读取上下文，但只有到达轴的末尾才能知道上下文的最终值。
当每个归约环节都能\emph{摆脱}其上下文时，就存在单遍算法：将它重写为上下文与不读取上下文的归约的函数。
本节给出实现这种分离的规则、应用规则的算法，以及由此得到的结论。

\subsection{Freeing a Reduction from Its Contexts}
\label{sec:freeing}

若操作 $\varphi$ 满足 $\varphi(a\oplus b)=\varphi(a)\oplus\varphi(b)$，则称其与归约 $\oplus$ \emph{可交换}。
\Cref{tab:rules} 列出 \sys 能够移到归约外部的操作；在附加条件成立且所有表达式有定义时，每条规则都是 $\R$ 上的恒等式。
求和对加法满足分配律，并允许提取与 $j$ 无关的因子。
最大值归约允许提取加法项、已证明非负的因子或非递减函数。
依赖上下文的项取偶次幂或绝对值时，可以转化为两个极值的较大者。

\begin{table}[t]
\footnotesize
\caption{Operations that commute with a reduction over $j$; $a$, $c$ and $y$ do not depend on $j$. Each rule is an identity under its side condition, which \sys proves syntactically before applying it.}
\label{tab:rules}
\begin{tabular}{@{}llp{0.42\linewidth}@{}}
\toprule
\textbf{Reduction} & \textbf{Rewrite} & \textbf{Side condition}\\
\midrule
$\sum_j$ & $\sum_j (u_j+v_j)=\sum_j u_j+\sum_j v_j$ & none\\
$\sum_j$ & $\sum_j a\,u_j=a\sum_j u_j$, \ $\sum_j e^{u_j+y}=e^{y}\sum_j e^{u_j}$ & none\\
$\max_j$ & $\max_j (u_j+a)=\max_j u_j+a$ & none\\
$\max_j$ & $\max_j c\,u_j=c\max_j u_j$ & $c\ge 0$ proven\\
$\max_j$ & $\max_j h(u_j)=h(\max_j u_j)$ & $h$ non-decreasing on the relevant domain\\
$\max_j$ & $\max_j |u_j|^{2k}=\max(\max_j u_j,\ \max_j -u_j)^{2k}$ & integer $k\ge1$\\
$\min_j$ & dual to $\max_j$ & dual\\
\bottomrule
\end{tabular}
\end{table}

下面的分离将归约体拆为 $r$ 个贡献项。
每个 $g_k$ 包含依赖输入数据的部分，$\varphi_{k,\ctx}$ 则承载该项对上下文的依赖。
重组时将上下文视为固定参数；因此，在尚不知道其值时，就可以先累积所得摘要。

\begin{lemma}[Regrouping]
\label{lem:regroup}
设 $f(x,\ctx)=\bigoplus_{k=1}^{r}\varphi_{k,\ctx}(g_k(x))$，其中所有 $g_k$ 都不读取上下文，且每个 $\varphi_{k,\ctx}$ 都与 $\oplus$ 可交换。
则对任意非空分区 $S$，有 $\bigoplus_{j\in S} f(x_j,\ctx)=\bigoplus_{k=1}^{r}\varphi_{k,\ctx}\bigl(\bigoplus_{j\in S} g_k(x_j)\bigr)$。
\end{lemma}

\begin{proof}
利用 $\oplus$ 的结合律和交换律重排归约，再将每个 $\varphi_{k,\ctx}$ 提取出来。
\end{proof}

所得元组无需知道最终上下文即可计算。上下文仅在从元组重构所需成员时使用。

\subsection{The Judgement}
\label{sec:judge}

参考原子是无上下文归约 $q(S)=\bigoplus_{j\in S}g(x_j)$，由其归约操作及规范化的归约体标识。
归约体允许为常量：$|S|=\sum_{j\in S}1$ 就是以加法为操作、以 $1$ 为归约体的原子。
收集原子前，按照部分语义，将所有作为数值读取的归约轴长度规范化为这一求和。
从求和中提取常量项也可能产生同一个原子。
这两种来源进入同一集合，并通过规范形式相等进行共享；只有所需表达式使用某个原子时，才将其加入集合。

\begingroup
\setlength{\intextsep}{2pt}
\reservewrapspace{16\baselineskip}
\begin{wrapfigure}{r}{0.48\textwidth}
\begin{minipage}{\linewidth}
\begin{algorithm}[H]
\small
\caption{\textsc{Judge}$(\mathit{chain}, \mathit{req})$}
\label{alg:judge}
\KwOut{reference atoms $\mathcal Q$, or failure}
Normalise extent reads in $\mathit{chain}$\;
$W\leftarrow \mathit{req}$;\ $\mathcal Q\leftarrow\emptyset$\;
\While{$W$ not empty}{
  $m\leftarrow$ pop $W$\;
  $e\leftarrow$ free contexts from body of $m$\;
  push unvisited dependencies of $e$ onto $W$\;
  \ForEach{atom $q$ in $e$}{
    \lIf{$q$ reads a context}{\Return \textsc{Stuck}$(m,q)$}
    $\mathcal Q\leftarrow\mathcal Q\cup\{q\}$ (canonical form)\;
  }
}
\Return \textsc{Fusible}$(\mathcal Q)$\;
\end{algorithm}
\end{minipage}
\end{wrapfigure}
\Cref{alg:judge} 访问每个所需成员及其依赖，每个成员只处理一次。
算法将上下文替换为沿归约轴保持不变的符号，通过 \cref{tab:rules} 重写归约体，并将剩余原子收集到 $\mathcal Q$ 中。
仍然读取上下文的原子称为\emph{受阻}：这些规则无法将它与上下文分离。
否则，收集到的原子构成参考状态元组，每个不同的原子对应一个状态。

这一构造给出有限摘要，其分量可以独立合并，其值能够确定所需输出。
它的大小界定了这一构造的状态数；状态选择将该上界用作默认状态数限制。
它也确定了验证所选表示时使用的无上下文量。
这两种作用需要区分：这一构造的存在性并不意味着每种保留原始结果的表示都能满足同一限制。

\finishwrapspace
\endgroup

\begin{theorem}[Existence]
\label{thm:existence}
设 \textsc{Judge} 返回 \textsc{Fusible}$(\mathcal Q)$，其中
\[
\mathcal Q=\{q_1,\dots,q_k\},\qquad
q_i(S)=\bigoplus_{j\in S}^{(i)}g_i(x_j).
\]
如果收集到的归约体在声明的元素域上都有定义，则该链存在至多使用 $|\mathcal Q|$ 个状态的单遍归约器。
其状态元组为 $(q_1(S),\dots,q_k(S))$，合并逐分量进行，收尾计算在源输出的定义域上重构所需输出。
\end{theorem}

\begin{proof}[Proof sketch]
轴长度的规范化保持部分语义；在已证明的附加条件下，每条代数重写在源定义域上都是 $\R$ 上的恒等式。
由 \cref{lem:regroup} 并按生产者顺序归纳可知，在源计算有定义的地方，每个所需成员都具有以收集到的原子表达的重构表达式。
对于每个原子，都有 $q_i(S\dunion T)=q_i(S)\oplus_i q_i(T)$。
因此，该元组可以逐分量合并，其叶子为 $(g_1(x_j),\dots,g_k(x_j))$。
定义域假设保证这些状态在合法非空分区上都有定义；在根节点应用重构表达式，即得到所需输出。
每个原子对应一个坐标，从而得到状态数上界。
\end{proof}

该判定识别能够通过这些规则分离上下文的归约链。
当剩余归约仍依赖某个上下文时，分析停止，不产生构造。
更多恒等式或更强的输入关系可能在这一已识别类别之外进一步建立可归约性。

这一构造是代数参考，不一定是适合实际执行的数值形式。
将上下文从归约体中分离会改变中间量：原本相消的项可能变成很大的独立求和，有界表达式可能变成无界因子的乘积。
保留源归约结果能够保留其中心化和缩放，而辅助状态提供合并所需的信息。
因此，下一步是推导实际维护的表示，并证明它保留了该构造所需的信息。
